% =====================================================================
% SECCION 8 -- DESARROLLO Y ENTRENAMIENTO DEL MODELO (esqueleto)
% Responsable: no corresponde al frente etico-legal.
%
% NOTA PARA EL EQUIPO: las decisiones tecnicas que se documenten aqui son
% auditadas en la Seccion \ref{sec:auditoria} del marco etico-legal.
% Cualquier cambio de arquitectura, de dataset o de politica de retencion
% debe comunicarse al frente etico-legal, porque cambia la auditoria.
%
% Los marcadores \pendiente{PLACEHOLDER}{...} que este archivo tenia se
% retiraron del cuerpo. El encargo consta en PENDIENTES.md, seccion (b).
% =====================================================================
% 14/09/2026: el aviso "Seccion en elaboracion ... a cargo de otro
% integrante" se sustituyo por la frase neutra comun a todos los
% esqueletos. Reestructuracion al indice del modelo del curso.
% Antes metodologia.tex. Conserva "Datos y entrenamiento" y "Metricas de
% evaluacion" sin cambiar sus palabras; "Arquitectura propuesta" paso a
% secciones/diseno_sistema.tex (Seccion 6).
\section{Desarrollo y Entrenamiento del Modelo}
\label{sec:desarrollo}

Esta sección se completará en la versión final del informe.

\subsection{Datos y entrenamiento}
\label{sec:datos-entrenamiento}

Este apartado documentará los conjuntos de datos utilizados, la
estrategia de entrenamiento y la de adaptación de dominio. La
verificación de las licencias de los conjuntos públicos, que la
Sección~\ref{sec:licencias} deja pendiente, no se ha completado todavía
y depende de que el equipo documente aquí qué conjuntos emplea y bajo qué
licencia se distribuye cada uno.

\subsection{Métricas de evaluación}
\label{sec:metricas}

Este apartado justificará las métricas seleccionadas (MAE, MSE, nAP u
otras). La Sección~\ref{sec:ant-p2pnet} señala una restricción que
conviene tener presente al elegirlas: los compromisos de exactitud deben
formularse a nivel de zona agregada y no de posición individual.
